\section{Exact optimization of square and single-rung rings}
\label{sec:four-coefficient-rings}

A square tile has two left vertices and two right vertices, with all four edges between the two pairs.
A single-rung tile has two vertices joined by one edge; these vertices serve as both its left and right pair.
Join the right pair of each tile to the next left pair by two matching edges, cyclically.
Each interface matching may be straight or crossed.
Let $a$ count square tiles and $s$ count single-rung tiles.
Put $W=2a+s$ and $m=a+s$.
The resulting cubic graph has rank $r=W+1$.
We restrict to $W\ge7$, so the small one-tile and two-tile degeneracies do not arise.

\begin{proposition}[Ordering the single-rung tiles]
\label{prop:four-coefficient-ring-order}
Suppose $a\ge1$.
Among all cyclic orders of these tiles, the maximum cycle excess over $2^W$ is
\[
 E(a,s)=2^a\left(2a+4s+\binom s2+\frac12\right)
        -3a-4s+\binom s2.
\]
It is attained when all single-rung tiles are consecutive.
The count is independent of the interface twists.
\end{proposition}

\begin{proof}
Assign weight $v_i=1/2$ to a square tile and $v_i=1$ to a single-rung tile.
For ordered tile positions $i,j$, let $v[i,j)$ be the product along the forward cyclic interval from $i$ through the position before $j$.
When $i=j$, use the empty interval and weight one.
Define
\[
 Q=\sum_{i,j=1}^{m}v[i,j).
\]

Cut parity makes every interface count odd, or makes all of them even.
The odd case uses one edge at each interface.
There are $2^m$ interface choices and two local paths in each square, hence exactly $2^{m+a}=2^W$ cycles.
A local path in a single-rung tile may have length zero when both interfaces choose the same vertex.

An even-sector cycle with no used interface is one of the $a$ squares.
If it uses every interface twice, each square has two through matchings.
Exactly half their $2^a$ choices make one circuit, giving $2^{a-1}$ cycles.
This remains true with any interface twists.

Otherwise the used interfaces form one nonempty proper cyclic interval.
Its two end tiles supply caps and its inner tiles supply through pairs.
A square has two choices in either role, and a single-rung tile has one.
The count is therefore $2^a$ times the product of the weights of the omitted tiles.
The omitted intervals have lengths zero through $m-2$.
The length-$m-1$ intervals, which are included in $Q$, have total weight
\[
 2^{-a}\sum_i v_i^{-1}=2^{-a}(2a+s)=2^{-a}W.
\]
Thus the cycle excess is exactly
\[
 a+2^{a-1}+2^aQ-W
 =2^a(Q+1/2)-m.
\]

Decompose $Q$ by the types of its two endpoints.
The square--square contribution is $2a(1-2^{-a})$.
For each single-rung tile, the two orders of mixed endpoints contribute
$2(1-2^{-a})$ and $1-2^{-a}$.
The equal single-rung endpoints contribute $s$.
For two distinct single-rung positions, their two cyclic interval weights are $z$ and $2^{-a}/z$, with $2^{-a}\le z\le1$.
Convexity gives
\[
 z+2^{-a}/z\le1+2^{-a}.
\]
Equality holds when one of those intervals contains no square.
All these equalities hold together when the single-rung tiles are consecutive.
Consequently
\[
 Q\le2a+4s+\binom s2
       +2^{-a}\left(-2a-3s+\binom s2\right).
\]
Substitution proves the formula and attainment.
\end{proof}

\begin{theorem}[The exact ring maximum at fixed rank]
\label{thm:four-coefficient-ring-maximum}
For $r\ge8$, every square and single-rung ring of rank $r$ has at most $H(r)$ cycles, where
\[
 \begin{aligned}
 H(2q+1)&=2^{2q}+(2q+1/2)2^q-3q,\\
 H(2q+2)&=2^{2q+1}+(2q+9/2)2^q-3q-4.
 \end{aligned}
\]
The maxima use no single-rung tile when $r$ is odd, and one when $r$ is even.
These are the Horak--Kim rings.
\end{theorem}

\begin{proof}
Fix $W=r-1\ge7$.
For $a\ge2$, compare $E(a,s)$ with $E(a-1,s+2)$, keeping $2a+s=W$.
Direct subtraction gives
\[
 E(a-1,s+2)-E(a,s)
 =-2^{a-1}\left(W+\frac{s^2+s}{2}-\frac{13}{2}\right)+2s-4<0.
\]
For $s=0,1$, the last inequality is immediate.
For $s\ge2$, the parenthesis exceeds $2s-4$ when $W\ge7$.
Thus among rings with at least one square, the maximum has $s=0$ or $s=1$ according to the parity of $W$.
Proposition~\ref{prop:four-coefficient-ring-order} gives the displayed formulas.

If $a=0$, every proper interval of used interfaces has one completion.
There are $W(W-1)$ such intervals.
The all-used case contributes at most one additional cycle.
Thus its excess is at most $W(W-1)+1$.
The stated maximum exceeds this quantity at $W=7,8$.
Its excess more than doubles when $W$ increases by two, whereas
\[
 (W+2)(W+1)+1<2\bigl(W(W-1)+1\bigr)\qquad(W\ge7).
\]
Induction on each parity proves the comparison for all $W\ge7$.
\end{proof}

\begin{remark}[A cap formula for open chains]
For an open chain with the same tile types, the through count is $D=2^a$.
Its disjoint left-cap/right-cap count is
\[
 C=2^a\sum_{i<j}2^{-a(i,j)},
\]
where $a(i,j)$ counts square tiles strictly between positions $i,j$.
The turn tiles contribute their local cap counts; all tiles before the first turn and after the second turn contribute through counts.
This gives the formula directly.
In particular,
\[
 C\le2^a\left(2a-4+2^{2-a}+4s+\binom s2\right).
\]
The square--square sum is exact, each single-rung tile contributes at most four against all squares, and each pair of single-rung tiles contributes at most one.
This sharpens the earlier bound to leading coefficient one when the number of single-rung tiles is minimal.
\end{remark}

\paragraph{Status.}
These formulas are new deductions from the previously audited saturated-ladder classification.
They have passed independent line audits by the root, path-stability, and random-host agents.
The companion checker verifies 527520 even states in 456 rings, 86 open chains, and 62491 exact optimization inequalities.
The proofs remain separate from the stable proof packet.
